\documentclass[12pt,a4paper]{article}
\usepackage[margin=2.2cm]{geometry}
\usepackage{fontspec}
\setmainfont{Times New Roman}
\usepackage{setspace}
\singlespacing
\usepackage{titlesec}
\titleformat{\section}{\normalfont\large\bfseries}{}{0em}{}
\titleformat{\subsection}{\normalfont\normalsize\bfseries}{}{0em}{}
\titlespacing{\section}{0pt}{8pt}{3pt}
\titlespacing{\subsection}{0pt}{6pt}{1pt}
\usepackage{hyperref}
\hypersetup{hidelinks}
\setlength{\parindent}{0pt}
\setlength{\parskip}{4pt}
\newenvironment{refs}{\begin{list}{}{\setlength{\leftmargin}{1.5em}\setlength{\itemindent}{-1.5em}\setlength{\itemsep}{0pt}\setlength{\parsep}{0pt}}\small}{\end{list}}

\begin{document}

{\large\bfseries BE451 Degree Project in Finance --- Thesis Synopsis, Fall 2026}

\textbf{Students:} [Full name (student ID)], [Full name (student ID)] \hfill \textbf{Preferred tutor:} [name / no preference]

\vspace{2pt}\hrule

\section{Topic 1: Investment and Payouts under a Threatened Profit Cap: Sweden's Private Welfare Firms}

\subsection{(i) Tentative research question}

\emph{How do firms change their investment and payouts when the government threatens to cap the profits they can earn on capital kept in the firm, and does the change reverse when the threat is withdrawn?}

Firms often delay investment when policy is uncertain, because waiting has option value (Julio and Yook, 2012; Jens, 2017). A threatened profit cap adds a second channel. If the return on capital kept in the firm may soon be limited by law, that capital is worth less to the owner. Owners then have a reason to pay out more while the cap is pending, and to stop when it is dropped. The two channels give different predictions: investment should fall under both, but payouts should rise only under the second. A third possibility is that owners never took the threat seriously, so nothing changes, which would show how little a politically weak threat costs.

Sweden's private welfare sector is a good place to test this. Every private provider of tax-financed schooling, social care and personal assistance faced a concrete proposal to cap operating profit at the government borrowing rate plus seven percentage points on operating capital. The proposal was published on 8 November 2016 (SOU 2016:78), became a government bill on 20 March 2018 (prop. 2017/18:159), and was rejected by the Riksdag in June 2018. Private health care was left out of the bill, so it is a natural comparison group with the same public funding. The threat returned for schools: an inquiry on schools without profit motive was appointed on 30 June 2022, and a bill of 25 June 2026 (prop. 2025/26:292) proposes a ban on value transfers from independent schools in three situations, to take effect on 1 July 2027. Welfare profits are also an issue in the election of 13 September 2026. We have found no study of how the firms responded.

\subsection{(ii) Potential data sources and empirical methods}

\textbf{Data.} Serrano (Swedish House of Finance) holds the annual accounts of all Swedish limited companies. We use 2012--2024: investment, dividends, group contributions, cash, debt, employees, entry, exit and group parent. Treated firms are private firms in education (SNI 85) and social care (SNI 87--88). The placebo group is private health care (SNI 86), which the bill excluded. Controls are matched private service firms. Payouts are measured at group level, since many welfare groups pay out through group contributions to holding companies, and we hand-code the owner type of the largest groups. The four listed welfare firms (Attendo, Humana, AcadeMedia, Ambea) allow a small stock-market event study using FinBas or LSEG. Serrano and FinBas are course resources [access verified on (date); about (N) treated firms].

\textbf{Method.} A difference-in-differences event study around the three dates in 2016--2018, with firm and industry-by-year fixed effects, comparing treated firms with the health-care placebo and with matched service firms. Outcomes: investment over assets, dividends and group contributions over equity, cash, employment growth, exit and ownership change. Mechanism tests: firms with pre-threat profits above versus below the proposed cap, and owner type. Robustness: placebo dates before 2016, excluding asylum accommodation (SNI 87.9), alternative matching. The 2022--2026 school episode is a second test, and the 2026 bill is a direct test of the payout channel because it bans payouts rather than capping profits. Main threats: the owner-screening rules in force from 1 January 2019 (main window ends in 2018), the interest-rate shock in the second episode (matching and the placebo group), and the small listed sample (illustrative only). Plan: 2014--2020 investment and payout results first, then mechanisms and the second episode.

\subsection{(iii) Short literature review}

\textbf{Julio and Yook (2012, Journal of Finance)} show across 48 countries that firms cut investment by 4.8 percent on average in election years and catch up afterwards. We use their event-window design on one dated policy threat inside one country, and add payouts and ownership changes as outcomes. \textbf{Jens (2017, Journal of Financial Economics)} shows the effect is causal using US gubernatorial elections and placebo tests; our placebo is the sector that the bill excluded. \textbf{Kelly, P\'astor, and Veronesi (2016, Journal of Finance)} show that political uncertainty is priced in options, and that the price depends on how likely the policy change is; we use this for the payout prediction and the listed-firm event study. \textbf{Bonaime, Gulen, and Ion (2018, Journal of Financial Economics)} show that policy uncertainty reduces merger activity; we use it for the ownership-change outcome. Our contribution is to test the investment channel and the payout channel on the full population of affected firms, with a placebo written into the bill itself, which election-based studies cannot do.

\begin{refs}
\item Bonaime, A., H. Gulen, and M. Ion, 2018, Does policy uncertainty affect mergers and acquisitions?, \emph{Journal of Financial Economics} 129, 531--558.
\item Jens, C. E., 2017, Political uncertainty and investment: Causal evidence from U.S. gubernatorial elections, \emph{Journal of Financial Economics} 124, 563--579.
\item Julio, B., and Y. Yook, 2012, Political uncertainty and corporate investment cycles, \emph{Journal of Finance} 67, 45--83.
\item Kelly, B., \v{L}. P\'astor, and P. Veronesi, 2016, The price of political uncertainty: Theory and evidence from the option market, \emph{Journal of Finance} 71, 2417--2480.
\item SOU 2016:78, Ordning och reda i v\"alf\"arden; prop. 2017/18:159; dir. 2022:102, En skola utan vinstintresse; prop. 2025/26:292, Sk\"arpta villkor f\"or friskolesektorn.
\end{refs}

\section{Topic 2: Lower Tax, Smaller Tax Shield: How Sweden's 2019 Corporate Tax Reform Changed the Leverage of Private Firms}

\subsection{(i) Tentative research question}

\emph{How did Sweden's 2019 corporate tax reform change the leverage of private firms, and was the response stronger among firms that risked losing part of their interest deductions?}

The reform did two things at once. The corporate tax rate fell from 22 to 21.4 percent in 2019 and to 20.6 percent in 2021, which makes the tax shield on debt less valuable for every firm. At the same time, net interest expense became deductible only up to 30 percent of taxable EBITDA, with a simplification rule that allows SEK 5 million of net interest regardless, shared within a group of associated companies. Standard trade-off theory predicts less debt from both changes. Ivanov, Pettit, and Whited (2026) find the opposite for small US private firms: after tax cuts, cash flows improve, default risk falls and leverage rises. Sanati and Beyhaghi (2025) find that US firms which lost part of their interest deduction in 2017 reduced both debt and investment. Sweden lets us study both channels in one reform and in the full population of limited companies. A zero response on both margins would also be informative: Swedish private-firm leverage would then not be tax-driven.

\subsection{(ii) Potential data sources and empirical methods}

\textbf{Data.} Serrano, all Swedish limited companies 2013--2024: debt, interest expense, tax expense, EBITDA, investment, dividends and group parent. Serrano is accounting data, not tax returns, so we cannot see exactly which firms were limited. We therefore build a \emph{predicted exposure} measure from 2017--2018 accounts: net interest expense relative to 30 percent of accounting EBITDA and to the SEK 5 million allowance, aggregated within groups. High-exposure firms are those whose pre-reform interest would have exceeded the new limit. The government's review of the rules (SOU 2024:37, pp.~78--82) reports from tax data that about 7,100 companies had deductions limited in 2019, that half of the limited companies were in real estate, and that the number of firms with net interest close to 30 percent of EBITDA rose from about 3,700 a year before the reform to about 6,100 after it. Serrano is a course resource [access verified on (date); about (N) firm-years].

\textbf{Method.} Main test: a difference-in-differences event study of leverage, new borrowing, interest expense over debt and investment on predicted exposure interacted with year indicators, with firm and industry-by-year fixed effects, and matched low-exposure firms as controls. Second test: the rate cut, using pre-reform tax payments over assets as the dose, so that firms that paid more tax before the reform gained more from the cut. Threats and how we handle them: exposure is predicted, not observed (we show a narrow and a broad definition, and check that the bunching at 30 percent seen in tax data also appears in accounting data); groups can shift interest between companies (exposure is measured at group level); the 2022 rate increases (industry-by-year effects and a placebo in 2016--2018); other 2019 rules aimed at multinationals (excluded). Plan: the exposure test first, then the rate cut, then real-estate heterogeneity. Fallback: exposure test only.

\subsection{(iii) Short literature review}

\textbf{Ivanov, Pettit, and Whited (2026, Review of Economic Studies)} use US state tax changes and private-firm data and find that leverage rises after tax cuts, because lower taxes raise the return on capital and lower credit spreads. We test their prediction on the 2019 and 2021 rate steps. \textbf{Sanati and Beyhaghi (2025, Review of Financial Studies)} use the USD 25 million exemption in the 2017 US reform and find that firms which lost the tax benefit of debt cut debt and investment without replacing debt with equity. Our exposure design follows the same logic, but with predicted exposure rather than a sharp threshold, since Sweden's SEK 5 million rule is an alternative to the EBITDA rule, not a cut-off. \textbf{Heider and Ljungqvist (2015, Journal of Financial Economics)} find that leverage rises by about 0.4 percentage points per percentage point of higher tax, but does not fall when taxes are cut. This asymmetry is why a small response to the rate cut would not be surprising, and why the deductibility channel is our main test. \textbf{Brok (2024, Review of Finance; 2026, Journal of Financial and Quantitative Analysis)} shows that group structure and tax avoidance change how debt responds to taxes, which is why we measure exposure at group level. Prior work on the Swedish reform: SOU 2024:37 documents the limited firms and the bunching with tax data, and finds descriptively, without a control group, that investment in the most affected groups did not fall (pp.~96--97); a master's thesis at Uppsala University (Jasem and Sado, 2026) finds modest negative leverage effects among large exposed firms in a firm-level difference-in-differences for 2016--2023. Our contribution relative to these is to use the full population rather than large firms, to assign exposure at group level, to separate the rate-cut channel from the deductibility channel, and to test the investment result with matched controls.

\begin{refs}
\item Brok, P., 2024, Debt and taxes: The role of corporate group structure, \emph{Review of Finance} 28, 2051--2082.
\item Brok, P., 2026, Debt and taxes: The role of tax avoidance, \emph{Journal of Financial and Quantitative Analysis} 61, 2007--2032.
\item Heider, F., and A. Ljungqvist, 2015, As certain as debt and taxes: Estimating the tax sensitivity of leverage from state tax changes, \emph{Journal of Financial Economics} 118, 684--712.
\item Ivanov, I. T., L. Pettit, and T. M. Whited, 2026, Taxes depress corporate borrowing: Evidence from private firms, \emph{Review of Economic Studies} 93, 2715--2749.
\item Jasem, S., and G. Sado, 2026, Leverage under constraint: Firm responses to Sweden's ATAD-based interest limitation reform, Master's thesis, Uppsala University.
\item Sanati, A., and M. Beyhaghi, 2025, How does removing the tax benefits of debt affect firms? Evidence from the 2017 U.S. tax reform, \emph{Review of Financial Studies}, forthcoming, doi:10.1093/rfs/hhaf041.
\item SOU 2024:37, F\"orb\"attrade r\"anteavdragsregler f\"or f\"oretag, 31 May 2024.
\end{refs}

\end{document}
